\abstract{РЕФЕРАТ}

Объем работы равен \formbytotal{lastpage}{страниц}{е}{ам}{ам}. Работа содержит \formbytotal{figurecnt}{иллюстраци}{ю}{и}{й}, \formbytotal{tablecnt}{таблиц}{у}{ы}{}, \arabic{bibcount} библиографических источников и \formbytotal{числоПлакатов}{лист}{}{а}{ов} графического материала. Количество приложений – 2. Графический материал представлен в приложении А. Фрагменты исходного кода представлены в приложении Б.

Перечень ключевых слов: интеллектуальная система, компьютерное зрение, сверточные нейронные сети, YOLOv11, Flask, PostgreSQL, веб-приложение, диагностика заболеваний, сельскохозяйственные растения, томаты, база данных, модуль, компонент, пользовательский интерфейс, машинное обучение, обработка изображений, Python, Google Colab, PyTorch, информационные технологии.

Объектом разработки является интеллектуальная система для диагностики заболеваний сельскохозяйственных растений, в частности томатов, с использованием технологий компьютерного зрения и машинного обучения.

Целью выпускной квалификационной работы является повышение эффективности диагностики заболеваний сельскохозяйственных растений путем создания интеллектуальной системы, которая позволяет загружать изображения растений, выявлять заболевания и предоставлять рекомендации по их лечению и профилактике.

В процессе создания системы были выделены основные сущности путем проектирования базы данных, использованы модули и классы, обеспечивающие обработку изображений и взаимодействие с базой данных, разработан пользовательский интерфейс на основе HTML/CSS и фреймворка Flask, а также реализована модель машинного обучения YOLOv11 для анализа изображений растений и выявления заболеваний.

При разработке системы использовались фреймворк Flask для серверной части, PostgreSQL для хранения данных о болезнях и рекомендациях, а также Python и PyTorch для реализации и обучения модели машинного обучения. 

Разработанная система была успешно протестирована.

\selectlanguage{english}
\abstract{ABSTRACT}
  
The volume of work is \formbytotal{lastpage}{page}{}{s}{s}. The work contains \formbytotal{figurecnt}{illustration}{}{s}{s}, \formbytotal{tablecnt}{table}{}{s}{s}, \arabic{bibcount} bibliographic sources and \formbytotal{числоПлакатов}{sheet}{}{s}{s} of graphic material. The number of applications is 2. The graphic material is presented in annex A. The layout of the site, including the connection of components, is presented in annex B.

List of keywords: intelligent system, computer vision, convolutional neural networks, YOLOv11, Flask, PostgreSQL, web application, disease diagnostics, agricultural plants, tomatoes, database, module, component, user interface, machine learning, image processing, Python, Google Colab, PyTorch, information technology.

The object of development is an intelligent system for diagnosing diseases of agricultural plants, in particular tomatoes, using computer vision and machine learning technologies.

The goal of the final qualifying work is to improve the efficiency of diagnosing diseases of agricultural plants by creating an intelligent system that allows you to upload plant images, identify diseases and provide recommendations for their treatment and prevention.

In the process of creating the system, the main entities were identified by designing a database, modules and classes were used to process images and interact with the database, a user interface based on HTML / CSS and the Flask framework was developed, and a YOLOv11 machine learning model was implemented to analyze plant images and identify diseases.

The system was developed using the Flask framework for the server side, PostgreSQL for storing disease and recommendation data, and Python and PyTorch for implementing and training the machine learning model.

The developed system was successfully tested.

\selectlanguage{russian}
